\chapter{Discussion}
\label{ch:discussion}

Neither SNPbag nor LDpred2-auto achieved meaningful discriminative
performance on the simulated phenotype, with AUCs of $0.54$ and $0.51$
respectively. The remainder of this chapter examines the limitations that
contributed to this outcome and identifies directions for future work.

The most significant limitation is the restricted genomic coverage.
The model was trained on only $14000$ SNPs from a single chromosome,
a constraint imposed entirely by GPU memory. Processing the full
chromosome~22 ($\sim\!116000$ SNPs) would exceed the memory of a
single NVIDIA B200 GPU even at the smallest batch size, and processing
all $22$ autosomes would require fundamentally different parallelisation
strategies. At the same time, the pre-training scaling experiment showed
that validation loss decreased monotonically as the training set grew
from $1000$ to $100000$ individuals, with the largest relative
improvement occurring between $10000$ and $100000$. This suggests
that the model has not yet saturated and would benefit from both longer
sequences and more training individuals.

All experiments were conducted on synthetic data generated by the HAPNEST
framework, not on real genotype data from a clinical cohort. While HAPNEST
preserves realistic linkage disequilibrium patterns and allele frequency
spectra, synthetic data may not fully capture the complexity of real human
genomes, including population substructure, rare variants, and genotyping
errors. The phenotype simulation used in this thesis assumes a simple
additive logistic model with a known set of causal variants, which may be
more favourable to both methods than the complex genetic architectures
underlying real diseases. Whether the pre-trained representations would
transfer to real data remains an open question.

Several architectural and methodological choices could be explored in
future work. SNPbag masks $85\%$ of genotype tokens during pre-training,
far higher than the $15\%$ used in the original BERT model
\citep{Devlin2019}. This aggressive masking forces the encoder to
reconstruct most of the sequence from very limited context, which may
strengthen the learned representations but could also make the
pre-training task unnecessarily difficult. The effect of the masking rate
on downstream phenotype prediction has not been investigated. Similarly,
the GELU activation function was inherited from the original BERT
architecture without exploration of alternatives such as SiLU or SwiGLU,
which have shown improvements in large language models. The prediction
head is a simple MLP applied to mean-pooled encoder output, and
alternative aggregation or prediction strategies could also be
considered. Each of these would require minimal code changes and could be
evaluated as straightforward ablation studies.

Finally, it is important to consider whether potential gains justify the
computational cost. Larger models and longer sequences may improve
prediction, but if the improvement is marginal or inconsistent across
traits, the added environmental and financial cost may not be worthwhile.
The results of this thesis suggest that the transformer architecture is
viable for learning genomic representations, the pre-training task
achieved $90\%$ masked accuracy, but that the current bottleneck is
scale, not architectural design.